\documentclass[acmsmall,nonacm,screen]{acmart}

\usepackage{mathpartir}

\citestyle{acmauthoryear}
\settopmatter{printacmref=false}

\newcommand{\eqd}{\doteq}
\newcommand{\kw}[1]{\mathsf{#1}}
\newcommand{\nm}[1]{\mathit{#1}}
\newcommand{\chk}[4]{#1 \vdash #2 : #3 \eqd #4}
\newcommand{\wf}[1]{\vdash #1}
\newcommand{\consts}{\mathrm{consts}}
\newcommand{\view}[1]{#1^{T,C}}
\newcommand{\lean}[1]{\texttt{#1}}

\begin{document}

\title{Reading Laws Instead of Bodies: Audit Surfaces Relative to an Interface}

\author{Claude-Fable-5.2}
\authornote{An AI model. It wrote this paper, the proofs, and the artifact at the request of a
human who asked for a paper that builds on the first paper of this repository. At the time of
writing nobody has reviewed it.}

\begin{abstract}
The audit surface of a machine-checked theorem holds the body of every definition that its
statement mentions. For a theorem about an implementation, the reader therefore reads the
implementation. We let the reader name some definitions as subjects and some theorems as laws,
and we define the interface surface, in which the subjects have no body and the laws that the
proof uses stand as axioms. We prove that the statement is derivable from the interface surface
whenever the proof, apart from the proofs of the laws, never unfolds a subject. The reader then
reads laws instead of bodies, and the kernel has checked the laws against the bodies. The result
is mechanized in Lean~4 on top of the proofs of the paper that it builds on.
\end{abstract}

\keywords{proof assistants, abstraction, interfaces, definitions, proof libraries}

\authorsaddresses{}

\maketitle

\section{Introduction}

A kernel checks the proof of a theorem, and a human who wants to rely on the theorem reads its
\emph{audit surface}~\citep{tpp1}: the definitions that the statement unfolds to and the axioms
that the proof uses. The statement is derivable from the audit surface alone, so nothing else in
the library has to be read.

The audit surface contains the body of every definition that the statement mentions. This is
right when the definition is vocabulary, such as the notion of a sorted list. It is a burden when
the definition is an implementation that the theorem is about. In the measurement
of~\citet{tpp1}, the median audit surface has 217 lines of code for the theorems of Lean's
\lean{Init} and 1,134 lines for those of \lean{Std}, whose statements mention the operations of
hash maps and tree maps. That paper calls such definitions the \emph{subjects} of a theorem and
leaves open how a reader may leave their bodies unread.

This paper closes that point. What a client knows about an abstract type is its
laws~\citep{mitchell1988}. We apply the same idea to reading. The reader names a set of
definitions as subjects and a set of theorems as \emph{laws}. The \emph{interface surface} of a
theorem contains no body of a subject. In their place it contains the laws that the proof uses,
as axioms. The reader reads the statements of these laws, and the kernel has checked their
proofs against the bodies.

Our main contributions are as follows:
\begin{itemize}
\item The interface surface of a checked proof for given laws and subjects
  (Section~\ref{sec:iface}).
\item The theorem that the statement is derivable from the interface surface when the proof,
  with every theorem outside the laws inlined, unfolds no subject (Theorem~\ref{thm:iface}), and
  a description of what the surface contains (Proposition~\ref{prop:entries}).
\item A mechanization in Lean~4~\citep{moura2021} that imports the mechanization
  of~\citet{tpp1} and reuses its lemmas.
\end{itemize}

\paragraph{Status of the claims.}
Theorem~\ref{thm:iface}, Proposition~\ref{prop:entries}, and the two surfaces of the example in
Section~\ref{sec:example} are mechanized, in about 400 lines that use no axioms beyond the three
standard axioms of Lean. The command \lean{make} in the artifact checks them. The discussion in
Section~\ref{sec:reading} is informal. The paper contains no measurement.

\section{Background and Example}
\label{sec:example}

We use the library calculus of~\citet{tpp1}. Terms are untyped $\lambda$-terms with constants. A
library $L$ is a sequence of definitions $\kw{def}\;c := u$, axioms
$\kw{axiom}\;a : s \eqd t$, and theorems $\kw{theorem}\;h : s \eqd t := p$. A proof term $p$
establishes an equation in a library, written $\chk{L}{p}{s}{t}$. A proof refers to an axiom
with $\kw{ax}\;a$, to a theorem with $\kw{lem}\;h$, and to the body of a definition with
$\kw{delta}\;c$. A library is well formed, written $\wf{L}$, when names are fresh, definitions
are closed and not recursive, and every theorem is checked against the entries before it. A set
of names is closed for $L$ when it contains the constants of the body of each definition and of
the statement of each axiom that it names. The audit set of a proof is the least closed set that
contains the constants of the statement and the axioms that the proof uses, and the audit surface
consists of the definitions and axioms that the set names.

\begin{figure}
\[
\begin{array}{ll}
\kw{def} & \nm{id} := \lambda x.\,x \\
\kw{def} & \nm{twice} := \lambda f.\,\lambda x.\,f\;(f\;x) \\
\kw{theorem} & \nm{twice\_id} : \nm{twice}\;\nm{id} \eqd \nm{id} \\
\kw{theorem} & \nm{client} : \nm{twice}\;\nm{id}\;(\nm{twice}\;\nm{id}) \eqd \nm{twice}\;\nm{id}
\end{array}
\]
\Description{A library with two definitions and two theorems.}
\caption{A library with a subject, a law, and a client.}
\label{fig:example}
\end{figure}

Figure~\ref{fig:example} shows a library. The theorem $\nm{twice\_id}$ is proved by normalizing
both sides, which unfolds $\nm{twice}$ and $\nm{id}$. The proof of $\nm{client}$ is
\[
\kw{trans}\;(\kw{app}\;(\kw{lem}\;\nm{twice\_id})\;(\kw{refl}\;u))\;
  (\kw{trans}\;(\kw{app}\;(\kw{delta}\;\nm{id})\;(\kw{refl}\;u))\;(\kw{beta}\;x\;u))
\qquad\text{with } u = \nm{twice}\;\nm{id} .
\]
It rewrites the head $\nm{twice}\;\nm{id}$ to $\nm{id}$ with the theorem, unfolds $\nm{id}$, and
contracts the redex. It never unfolds $\nm{twice}$.

The audit surface of $\nm{client}$ consists of the definitions of $\nm{twice}$ and $\nm{id}$,
since the statement mentions both. A reader who takes $\nm{twice}$ as a given artifact would
rather read that $\nm{twice}\;\nm{id} \eqd \nm{id}$ holds of it, together with the definition of
$\nm{id}$.

\section{The Interface Surface}
\label{sec:iface}

Fix a set $T$ of theorem names, the \emph{laws}, and a set $C$ of names of definitions, the
\emph{subjects}.

The \emph{view} $\view{L}$ of a library $L$ is obtained in three steps. Remove the theorems
outside $T$. Replace each theorem $\kw{theorem}\;h : s \eqd t := q$ in $T$ by the axiom
$\kw{axiom}\;h : s \eqd t$. Remove the definitions of the subjects. In the view, a subject is a
constant without a definition, and a law is an assumption.

For a proof $p$, let $p^T$ be the proof in which every $\kw{lem}\;h$ with $h \notin T$ is
replaced by the proof of $h$ in $L$, for the theorems of $L$ from the newest to the oldest, and
every remaining $\kw{lem}\;h$ is read as $\kw{ax}\;h$.

\begin{definition}[Interface surface]
\label{def:iface}
Let $\chk{L}{p}{s}{t}$. The \emph{interface set} of $p$ is the least set that is closed for
$\view{L}$ and contains $\consts(s)$, $\consts(t)$, and the axioms that $p^T$ refers to. The
\emph{interface surface} of $p$ consists of the definitions and the axioms of $\view{L}$ whose
names are in the interface set, in the order of $L$.
\end{definition}

The axioms that $p^T$ refers to are the axioms that $p$ uses outside the proofs of the laws,
and the laws that it uses. Since $\view{L}$ is well formed when $L$ is, one pass over it computes
the interface set, by Proposition~3.3 of~\citet{tpp1}.

\begin{theorem}[Interface soundness]
\label{thm:iface}
Let $\wf{L}$ and $\chk{L}{p}{s}{t}$. Suppose that the proof obtained from $p$ by inlining the
theorems outside $T$ refers to the definition of no subject. Then the interface surface of $p$
is well formed and contains no theorem, and there is a proof $p'$ that establishes $s \eqd t$ in
it.
\end{theorem}

\begin{proof}[Proof outline]
Inline the theorems outside $T$ into $p$, from the newest to the oldest. As in the proof of
audit soundness~\citep{tpp1}, each step replaces $\kw{lem}\;h$ by a proof of the same equation.
The result checks in $L$ without the theorems outside $T$. Reading each remaining $\kw{lem}\;h$
as $\kw{ax}\;h$ gives the proof $p^T$, which checks in the library whose laws are axioms. By
assumption $p^T$ contains no $\kw{delta}\;c$ for a subject $c$, so it still checks when the
definitions of the subjects are removed, by the locality of checking. This library is $\view{L}$.
It is well formed and contains no theorem, and the interface set is closed for it and contains
the constants of the statement and the axioms that $p^T$ refers to. Definition elimination, the
second half of the proof of audit soundness, then gives a proof over the interface surface.
\end{proof}

\begin{proposition}[Contents of the surface]
\label{prop:entries}
The interface surface defines no subject. Each of its entries is a definition of $L$, an axiom of
$L$, or the axiom $\kw{axiom}\;h : s \eqd t$ for a theorem $\kw{theorem}\;h : s \eqd t := q$ of
$L$ with $h \in T$.
\end{proposition}

With no law and no subject, the hypothesis of Theorem~\ref{thm:iface} holds for every proof, and
the theorem states that the statement is derivable from a closed set of definitions and axioms,
as audit soundness does.

\paragraph{The example.}
Take $T = \{\nm{twice\_id}\}$ and $C = \{\nm{twice}\}$ in Figure~\ref{fig:example}. The proof of
$\nm{client}$ refers to the law and to the definition of $\nm{id}$, so the hypothesis of
Theorem~\ref{thm:iface} holds. The interface surface is
\[
\kw{def}\;\nm{id} := \lambda x.\,x, \qquad
\kw{axiom}\;\nm{twice\_id} : \nm{twice}\;\nm{id} \eqd \nm{id} .
\]
The body of $\nm{twice}$ is absent. A second proof of the same statement normalizes both sides.
It unfolds $\nm{twice}$, so Theorem~\ref{thm:iface} does not apply to it, and its audit surface
is the right thing to read.

\section{What the Reader Learns}
\label{sec:reading}

\emph{Two facts.} A reader of the interface surface learns that the statement follows from the
laws, the vocabulary, and the axioms that are listed. The reader also knows that the library is
well formed, so the kernel has checked each law against the bodies of the subjects. Together
these say that the statement holds of the subjects as they are in the library. The reader has
read no body of a subject.

\emph{The price.} The laws are now part of what is read. Their statements can mention further
definitions, which the interface set contains through its closure. The interface surface is
smaller than the audit surface when the laws that a proof uses are shorter than the bodies they
replace.

\emph{Choosing subjects.} Theorem~\ref{thm:iface} holds for every choice of $T$ and $C$ that
meets its hypothesis. Two choices deserve care. A definition whose body the proof unfolds
outside the laws cannot be a subject, and the hypothesis rules it out. A definition that an
axiom of the surface mentions can formally be a subject, and then the axiom assumes something
about a constant whose body the reader has not seen. The reader should keep such a definition as
vocabulary, as~\citet{tpp1} recommends.

\emph{Replacement.} A subject can be given a new body together with new proofs of its laws, and
its clients stay checked~\citep{tpp1}. The interface surface of a client mentions neither the
old body nor the old proofs, so we expect it to be the same before and after. We have not
mechanized this statement.

\section{Limits}

\emph{Explicit unfolding.} The hypothesis of Theorem~\ref{thm:iface} can be read off a proof
term, because our proofs record each use of a definition with $\kw{delta}$. In a proof assistant
such as Lean, unfolding happens inside the conversion check of the kernel, and a proof term does
not show which bodies it relies on. There the system has to enforce the hypothesis, with
definitions that are opaque or that a module exports without their bodies.

\emph{No measurement.} We do not know how much smaller interface surfaces are than audit surfaces
in a real library. The two medians quoted in the introduction only suggest that implementations
account for much of the reading.

\emph{The choice of laws.} The theorem takes the laws as given. It does not say whether a set of
laws describes a subject well, and a client can only be read through the laws that its proof
happens to use.

\emph{The calculus.} The calculus is untyped and equational, with the limits that~\citet{tpp1}
lists.

\section{Related Work}

\citet{mitchell1988} explain abstract types as existential types: a client is checked against
the operations and never against the representation. Our laws play the role of the operations'
specification, and Theorem~\ref{thm:iface} is the corresponding statement for reading.
\citet{leroy1994} distinguishes manifest from abstract components of a module, which is the
distinction between vocabulary and subjects. \citet{gratzer2022} give a type theory in which the
unfolding of a definition is controlled, which is one way to enforce our hypothesis.
\citet{pollack1998} separates checking a proof from understanding what it proves. The audit
surface~\citep{tpp1} makes the second task precise, and the interface surface adapts it to
theorems about implementations.

\section{Conclusion}

A theorem about an implementation can be read through the laws of the implementation. The
interface surface replaces the bodies of the subjects by the laws that the proof uses, and the
statement is derivable from it when the proof unfolds no subject outside the laws. The next steps
are a measurement in a library whose modules hide bodies, and the invariance of the interface
surface under the replacement of a subject.

\bibliographystyle{ACM-Reference-Format}
\bibliography{refs}

\end{document}
